\chapter{Dynamic Network Scanning and Topology Inference}
\label{chap:network_scanner}

\section{Real-Time Discovery Pipeline}
Static vulnerability databases fail to capture dynamic network reconfigurations. To ground the environmental assessment in physical reality, the HelpUS platform integrates an active network interrogation and topology discovery module.

The discovery subsystem interrogates target subnets across four distinct operational modes:
\begin{enumerate}
    \item \textbf{Local Host Probe:} Interrogates the operator or visitor environment via native browser capabilities, client platform hints, and local socket probes to establish identity and loopback boundary.
    \item \textbf{Subnet CIDR Range Sweep:} Scans arbitrary private subnets (e.g., $192.168.0.0/24$, $10.0.0.0/16$) via concurrent socket connections across standard enterprise service ports (HTTP 80/443, SSH 22, PostgreSQL 5432, RDP 3389, MySQL 3306).
    \item \textbf{Enterprise Presets:} Documented baseline topologies including PCI-segmented cardholder environments, DMZs, and network-isolated clinical intranets. These presets are controlled fixtures for development and demonstration, not independently validated production topologies.
    \item \textbf{Audit JSON Ingestion:} Direct consumption of Nmap, OpenVAS, or Qualys output artifacts for automated continuous integration.
\end{enumerate}

\section{Topology Inference and Attack Surface Mapping}
Once active hosts and open listener ports are discovered, the topology engine derives the reachability graph:
\begin{equation}
    \mathcal{G} = (\mathcal{V}_{\text{nodes}}, \mathcal{E}_{\text{edges}})
\end{equation}
where nodes $\mathcal{V}_{\text{nodes}}$ represent gateways, firewalls, and workloads, and edges $\mathcal{E}_{\text{edges}}$ represent verified routing paths.

\subsection{Exposure Classification Logic}
A host $h \in \mathcal{V}_{\text{nodes}}$ is classified based on routing and perimeter boundaries:
\begin{itemize}
    \item \textbf{Direct Internet Exposure ($I_{\text{exposed}} = \text{true}$):} Verified inbound path from an external network to a service bound on an Internet-routable address or explicitly published through a gateway/NAT. The CIDR $0.0.0.0/0$ denotes the IPv4 default route/address set and is not itself a host address.
    \item \textbf{Segmented Subnet ($N_{\text{segmented}} = \text{true}$):} Workload located behind a stateful gateway or VLAN tag, requiring intra-subnet lateral movement for exploitation.
    \item \textbf{Local Loopback ($I_{\text{exposed}} = \text{false}, N_{\text{segmented}} = \text{true}$):} Service bound strictly to loopback/IPC with no verified remote listener. This supports a local-reachability hypothesis only when vulnerability semantics are consistent with local exploitation.
\end{itemize}

\section{Throughput and Latency Performance}

The current implementation records engine-level timings, but the manuscript does not present them as a controlled speedup against historical human analysis because the quantities were collected under different protocols. A publication-quality latency benchmark should report the engine hardware/software environment, repeated runs, and distributional statistics (for example median and p95), and should compare equivalent end-to-end tasks.

\section{Next-Generation Firewall (NGFW) Telemetry Ingestion}
\label{sec:ngfw_telemetry}
While active port sweeping identifies live listener sockets, enterprise defense in depth relies fundamentally on stateful packet inspection, Intrusion Prevention Systems (IPS), and layer-7 application firewalls. To bridge passive logs with active prioritization, HelpUS defines an ingestion adapter for structured Next-Generation Firewall (NGFW) and IDS telemetry. Production validation against specific vendors must be documented separately with versioned artifacts before it is claimed as an empirical result.

Let $\mathcal{L} = \{\ell_1, \ell_2, \dots, \ell_k\}$ represent the stream of structured firewall event records. Each event $\ell_i$ is represented by the tuple:
\begin{equation}
    \ell_i = \langle \text{timestamp}, \text{IP}_{\text{src}}, \text{IP}_{\text{dst}}, \text{port}_{\text{dst}}, \text{proto}, \text{action}, \text{threat\_id}, \text{rule\_name} \rangle
\end{equation}
where $\text{action} \in \{\text{allow}, \text{deny}, \text{drop}, \text{reset}\}$.

The topology engine correlates vulnerability listener ports with directional firewall session statistics:
\begin{equation}
    \Psi(h, p) = \frac{\sum_{\ell \in \mathcal{L}} \mathbb{I}(\text{dst}=h \land \text{port}=p \land \text{action}=\text{allow})}{\sum_{\ell \in \mathcal{L}} \mathbb{I}(\text{dst}=h \land \text{port}=p)}
\end{equation}
If $\Psi(h, p)=0$ over a defined observation window, the result is stored as evidence that no allowed sessions were observed in that telemetry source during that period. It does \emph{not} by itself prove absence of reachability and does not directly transform $\text{MAV:Network}$ into $\text{MAV:Adjacent}$ or $\text{MAV:Local}$. MAV changes require corroborating topology and vulnerability-semantic evidence; incomplete observations trigger the fail-safe review policy.

